\chapter{Patterns: state colocation}
\label{ch:colocation}

\begin{aside}
The closer state lives to where it's used, the easier to
reason about. Hoist only when sharing demands it.
\end{aside}

\section{The rule}

State lives at the lowest scope where it's read. Two
widgets need it? Hoist to common parent. Five widgets in
disjoint subtrees? Hoist to root (store).

\section{Component-local}

A counter widget owns its count. Other widgets shouldn't
need to know:

\begin{lstlisting}
auto counter = [&] {
    auto count = signal(0);
    return column(
        text(std::to_string(ctx.get(count))),
        button("+", [&] {...})
    );
};
\end{lstlisting}

\section{Sibling sharing}

Two siblings share state via the parent:

\begin{lstlisting}
auto parent = [&] {
    auto shared = signal(0);
    return row(producer(shared), consumer(shared));
};
\end{lstlisting}

\section{Tree-wide}

A store at root scope; consumers anywhere read.

\section{The path test}

If a piece of state is passed through three or more layers
of components without being touched, lift to a store.

\section{Recap}

\begin{recap}
\begin{itemize}[leftmargin=*]
  \item Local first; hoist only when needed.
  \item Sharing siblings via common parent.
  \item Wide sharing via store.
  \item Prop drilling > 3 layers = lift.
\end{itemize}
\end{recap}

\section*{Workshop}

\begin{enumerate}[leftmargin=*]
  \item Find a signal in your app that's too high.
    Lower it.
  \item Find prop drilling. Lift to a store.
\end{enumerate}
